% Lösungen Einheit 2 von 3 – Rückwärts (zusammenbau v0.1)
\einheitenkopf{Einheit 2 von 3 · Rückwärts}

\erg{12}{a) nutzen – aus der Unabhängigkeit eine Gleichung gewinnen: $0{,}3 \cdot P(B) = 0{,}12$ \quad b) $35\,\%$ – bei Unabhängigkeit ist der Anteil in der Teilgruppe gleich dem Gesamtanteil. \quad c) $P(A) \cdot P(B) = (p + 0{,}2) \cdot 4p = p$ führt auf $4p^2 - 0{,}2p = 0$, also $p = 0$ oder $p = 0{,}05$; $p = 0$ entfällt (dann wäre $P(B) = 0$), also $p = 0{,}05$. \quad d) $1 - 2p = 2p^2 + (1 - 2p)^2$ (links $P_E(G)$: nach der 3 gewinnt nur die 3; rechts $P(G)$ als Pfadsumme); daraus $6p^2 - 2p = 0$, also $p = \frac{1}{3}$ ($p = 0$ entfällt). \quad e) $P(B) = \frac{0{,}1}{0{,}4} = 0{,}25$ \quad f) $0{,}3 = 0{,}2 + x - 0{,}2x$, also $0{,}1 = 0{,}8x$ und $x = 0{,}125$ \quad g) Mit $x = P(A)$: $P(\overline{B}) = 0{,}7 - x$; Unabhängigkeit liefert $x \cdot (0{,}7 - x) = 0{,}1225$, also $x^2 - 0{,}7x + 0{,}1225 = 0$ und $x = 0{,}35$ (doppelte Lösung).}
\erg{13}{a) Fehler: Mia bildet die Differenz der Randwahrscheinlichkeiten. Richtig: $P(A \cap \overline{B}) = P(A) \cdot P(\overline{B}) = 0{,}5 \cdot 0{,}7 = 0{,}35$ \quad b) $P(A)$ zerfällt in $P(A \cap B) + P(A \cap \overline{B})$; der erste Summand ist $P(A) \cdot P(B)$, also bleibt $P(A) - P(A) \cdot P(B) = P(A) \cdot (1 - P(B))$, und das ist $P(A) \cdot P(\overline{B})$. \quad c) $0{,}99 = 0{,}9 + x - 0{,}9x$, also $0{,}09 = 0{,}1x$ und $x = 0{,}9$ – der zweite Sensor muss ebenfalls mit $0{,}9$ erkennen.}
